% =====================================================================
%  KỊCH BẢN THUYẾT TRÌNH 20 PHÚT — bản in để cầm đọc
%
%  Biên dịch: pdflatex -interaction=nonstopmode kich-ban-thuyet-trinh-20-phut.tex
%             (CHẠY HAI LẦN cho mục lục đúng)
%
%  Nội dung lấy nguyên từ kich-ban-thuyet-trinh-20-phut.md, chỉ đổi cách
%  trình bày: chữ to hơn, lời NÓI tách hẳn khỏi chỉ dẫn sân khấu.
%
%  MỌI CON SỐ GÕ TAY từ ket-qua/*.csv. Chạy lại pipeline thì phải đối
%  chiếu lại — khác với slide và web (đọc thẳng từ pipeline).
% =====================================================================
\documentclass[12pt]{article}

\usepackage[utf8]{inputenc}
\usepackage[T5]{fontenc}
\usepackage[vietnamese]{babel}
\usepackage[a4paper,margin=2cm,bottom=2.2cm]{geometry}
\usepackage{amsmath}
\usepackage{booktabs}
\usepackage{array}
\usepackage{enumitem}
\usepackage{xcolor}
\usepackage{fancyhdr}
\usepackage{hyperref}

\hypersetup{unicode=true, hidelinks,
  pdftitle={Kich ban thuyet trinh 20 phut}}
\setlist{leftmargin=*, topsep=2pt, itemsep=2pt}
\emergencystretch=3em
\setlength{\parindent}{0pt}
\setlength{\parskip}{4pt}

\definecolor{xam}{gray}{0.45}
\definecolor{vach}{gray}{0.75}

\pagestyle{fancy}
\fancyhf{}
\fancyhead[L]{\small\color{xam}Kịch bản thuyết trình — 20 phút}
\fancyhead[R]{\small\color{xam}\thepage}
\renewcommand{\headrulewidth}{0.4pt}

% ---- Tiêu đề một slide: số, tên, thời lượng, ghi chú cắt/không cắt ----
% Tham số 4 để rỗng nếu slide bình thường.
\newcommand{\slidemuc}[4]{%
  \vspace{10pt}%
  {\color{vach}\hrule height 0.8pt}%
  \vspace{4pt}%
  {\large\bfseries Slide #1 — #2}%
  \hfill{\color{xam}\normalsize #3}%
  \ifx\relax#4\relax\else\\[2pt]{\bfseries\small #4}\fi%
  \vspace{4pt}\par%
}

% ---- Khối LỜI NÓI: gạch dọc bên trái, thụt vào, chữ to hơn ----
\newenvironment{noi}%
  {\vspace{3pt}\begin{list}{}{%
     \setlength{\leftmargin}{14pt}\setlength{\rightmargin}{0pt}%
     \setlength{\itemindent}{0pt}\setlength{\listparindent}{0pt}%
     \setlength{\parsep}{5pt}}\item[]%
   \color{black}}%
  {\end{list}\vspace{2pt}}

% ---- Chỉ dẫn sân khấu và câu chuyển ----
\newcommand{\chuyen}[1]{\vspace{2pt}{\color{xam}\small$\rightarrow$\ \emph{#1}}\par}
\newcommand{\chidan}[1]{{\color{xam}\small\emph{#1}}\par}

\begin{document}

\begin{center}
  {\LARGE\bfseries Kịch bản thuyết trình}\\[4pt]
  {\large Thuế GTGT 10\% xuống 8\% và giá bán lẻ --- Nhóm 5}\\[6pt]
  {\color{xam}20 phút, gồm demo trực tiếp. Bám theo 25 slide ở \texttt{/trinh-bay}.}
\end{center}

\vspace{6pt}

\begin{center}
\begin{tabular}{@{}lcc@{}}
\toprule
\textbf{Đoạn} & \textbf{Slide} & \textbf{Phút} \\
\midrule
Mở đầu và dữ liệu        & 1--7   & 4:05 \\
Đại lượng đo và vấn đề   & 8--10  & 2:40 \\
Phương pháp              & 11--19 & 5:40 \\
Kết quả                  & 20--23 & 3:40 \\
\textbf{Demo trực tiếp}  & ---    & \textbf{2:30} \\
Hạn chế và kết luận      & 24--25 & 1:15 \\
\midrule
& & \textbf{19:50} \\
\bottomrule
\end{tabular}
\end{center}

\vspace{4pt}

\textbf{Nếu bị hối, cắt theo thứ tự:} slide 5 $\to$ slide 13 $\to$ slide 4. Ba slide
này là chi tiết kỹ thuật, bỏ đi không ai hỏi. Cắt cả ba thì dôi ra 1 phút 20.

\textbf{Đừng bao giờ cắt} slide 9, 10, 21, 23 --- cắt là mất mạch lập luận.

\vspace{4pt}

\chidan{Mọi con số trong kịch bản được gõ tay từ \texttt{ket-qua/*.csv} và đã đối
chiếu 22 trên 22. Chạy lại pipeline thì phải kiểm lại.}

% =====================================================================
\section*{Phần 1 --- Mở đầu và dữ liệu}

\slidemuc{1}{Bìa}{20s}{\relax}
\begin{noi}
Chào thầy và các bạn. Nhóm 5 làm đề tài: việc giảm thuế giá trị gia tăng từ 10\%
xuống 8\% có thật sự làm giảm giá bán lẻ mà người mua phải trả không.
\end{noi}
\chuyen{Vì sao câu hỏi này không hiển nhiên như nghe qua.}

\slidemuc{2}{Đặt vấn đề}{50s}{\relax}
\begin{noi}
Nhà nước giảm thuế để hỗ trợ người tiêu dùng. Nhưng nhà nước chỉ quyết định thuế
suất, còn cửa hàng mới là bên quyết định giá niêm yết. Hai quyết định đó độc lập
với nhau.

Nên chuyện thuế giảm mà giá có giảm theo hay không là một câu hỏi \textbf{thực
nghiệm}, phải đi đo, chứ không suy ra được.
\end{noi}
\chuyen{Nhóm đo bằng dữ liệu gì.}

\slidemuc{3}{Dữ liệu}{40s}{\relax}
\begin{noi}
Hóa đơn điện tử của một cửa hàng tiện lợi ở TP.HCM, ghi ở cấp từng dòng hàng. Mỗi
dòng có mã hàng, thuế suất được áp, và số tiền sau thuế.

Cấp dòng là điểm quan trọng: nó cho phép theo dõi \textbf{cùng một mã hàng} trước
và sau ngày chính sách, thay vì so hai giỏ hàng khác nhau.
\end{noi}
\chuyen{Dữ liệu thô phải qua vài bước lọc.}

\slidemuc{4}{Xử lý dữ liệu}{30s}{\itshape cắt được}
\begin{noi}
Mỗi bước lọc đều ghi rõ quy tắc và số dòng còn lại, để ai đọc cũng dựng lại được.
Từ 233 nghìn dòng thô còn lại \textbf{82 nghìn dòng hóa đơn}.

Lưu ý đây là số \textbf{dòng}, chưa phải số mặt hàng --- một chai dầu gội bán 300
lần thì là 300 dòng. Gộp lại theo mã hàng thì ra \textbf{2.218 mặt hàng}, và lát
nữa nhóm sẽ lọc tiếp còn \textbf{287 mặt hàng} đem so sánh.
\end{noi}
\chuyen{Trước khi chọn cửa sổ thời gian, nhóm rà độ phủ.}

\slidemuc{5}{Phân bổ theo thời gian}{25s}{\itshape cắt được}
\begin{noi}
Nhóm đếm số hóa đơn theo từng tháng để xem đoạn nào dữ liệu mỏng, rồi mới chọn cửa
sổ phân tích --- chứ không chọn trước rồi mới nhìn.
\end{noi}
\chuyen{Còn một thứ phải kiểm trước khi phân tích.}

\slidemuc{6}{Đánh giá sơ bộ}{40s}{\relax}
\begin{noi}
Nhóm đếm thuế suất ghi trên hóa đơn của từng mặt hàng, trước và sau ngày chính
sách. Đây là chỗ phát hiện ra một chuyện quan trọng: \textbf{có những mặt hàng luật
cho giảm nhưng cửa hàng vẫn xuất 10\%.}

Nên phải tách bạch hai thứ: mặt hàng có \textbf{đủ điều kiện} giảm thuế theo luật,
và thuế suất cửa hàng \textbf{thực sự áp}. Hai cái này không trùng nhau.
\end{noi}
\chuyen{Cơ sở pháp lý là thứ tạo ra nhóm đối chứng.}

\slidemuc{7}{Cơ sở pháp lý}{40s}{\relax}
\begin{noi}
Nghị quyết 204 năm 2025, hiệu lực mùng 1 tháng 7. Điểm hay là nó \textbf{không giảm
cho tất cả}: rượu bia, thuốc lá, xăng dầu, viễn thông vẫn giữ 10\%.

Chính chỗ loại trừ đó cho nhóm mình một nhóm đối chứng có sẵn --- cùng cửa hàng,
cùng khoảng thời gian, chỉ khác ở chỗ được giảm thuế hay không. Đây là một
\textbf{thí nghiệm tự nhiên}.
\end{noi}
\chuyen{Vậy đo cái gì trên từng mặt hàng.}

% =====================================================================
\section*{Phần 2 --- Đại lượng đo và vấn đề}

\slidemuc{8}{Đại lượng cần đo}{55s}{\relax}
\begin{noi}
Mỗi mặt hàng cho đúng một con số: giá sau so với giá trước, tính bằng log rồi nhân
100.

Dùng log chứ không lấy hiệu giá thô vì hai lý do. Một là giá bán lẻ lệch phải rất
mạnh --- vài mặt hàng đắt gấp trăm lần mặt hàng rẻ, nếu lấy hiệu thì mấy món đắt
chi phối hết. Hai là ta quan tâm \textbf{thay đổi bao nhiêu phần trăm}, chứ không
phải bao nhiêu đồng.

Nhân 100 chỉ để con số dễ đọc: một điểm log xấp xỉ một phần trăm.
\end{noi}
\chuyen{Đo xong thì ra thế này.}

\slidemuc{9}{Kết quả đo được}{55s}{KHÔNG CẮT}
\begin{noi}
Mẫu so sánh gồm \textbf{287 mặt hàng} --- phần lớn kho hàng vốn đã chịu 8\% từ
trước nên không có gì để so, nhóm chỉ giữ lại những mặt hàng đổi thuế suất đúng dịp
này.

Đây là chỗ đầu tiên gây bất ngờ. Nhóm được giảm thuế: 155 mặt hàng, giá trung bình
\textbf{tăng 0,62 điểm}. Nhóm không được giảm: 132 mặt hàng, \textbf{tăng 1,02
điểm}.

Cả hai nhóm đều \textbf{tăng} giá. Không nhóm nào giảm.

Nhưng nhóm được giảm thuế tăng \textbf{ít hơn} --- chênh lệch là 0,40 điểm.
\end{noi}
\chuyen{Câu hỏi tiếp theo là 0,40 điểm đó có phải do chính sách không.}

\slidemuc{10}{Vì sao chưa dùng thẳng được}{50s}{KHÔNG CẮT}
\begin{noi}
Chưa. Vì hai nhóm này khác nhau sẵn từ đầu --- một bên là rượu bia thuốc lá, một
bên là hàng tiêu dùng thường ngày. Chúng vốn đã có nhịp đổi giá khác nhau, kể cả
khi \textbf{không có chính sách nào cả}.

Nên 0,40 điểm này đang trộn hai thứ: phần do chính sách, và phần do hai nhóm vốn
khác nhau. Việc còn lại của bài là tách hai phần đó ra.
\end{noi}
\chuyen{Nhóm dùng khung nào để tách.}

% =====================================================================
\section*{Phần 3 --- Phương pháp}

\slidemuc{11}{Khung phương pháp}{45s}{\relax}
\begin{noi}
Ý tưởng: lấy mức đổi giá của nhóm được giảm thuế, trừ đi mức đổi giá của nhóm đối
chứng. Phần còn lại là ước lượng tác động.

Cách này đứng được là nhờ một giả định: \textbf{nếu không có chính sách, hai nhóm
sẽ đổi giá theo cùng một xu hướng.} Lát nữa nhóm sẽ nói thẳng là giả định này có
vấn đề.
\end{noi}
\chuyen{Để nói rõ cái gì cần kiểm soát, nhóm vẽ đồ thị nhân quả.}

\slidemuc{12}{Đồ thị nhân quả}{45s}{\relax}
\begin{noi}
Đây là bản đồ cái gì ảnh hưởng cái nào.

Điểm cần chú ý: \textbf{thuế suất cửa hàng thực áp nằm ở giữa} --- giữa việc mặt
hàng có đủ điều kiện giảm thuế và giá cuối cùng. Nó là \textbf{biến trung gian}.

Vì thế nhóm ước lượng theo \textbf{đủ điều kiện theo luật}, chứ không hồi quy theo
thuế suất thực áp. Kiểm soát biến trung gian là chặn mất đúng con đường mình cần đo.
\end{noi}
\chuyen{Từng biến trong sơ đồ được đo bằng gì.}

\slidemuc{13}{Gắn biến vào dữ liệu}{25s}{\itshape cắt được}
\begin{noi}
Mỗi biến trong khung ứng với một trường cụ thể trên hóa đơn --- không có biến nào
là suy diễn hay gán tay.
\end{noi}
\chuyen{Nhóm ước lượng bằng bốn cách.}

\slidemuc{14}{Bốn cách --- tổng quan}{35s}{\relax}
\begin{noi}
Bốn cách này \textbf{giống nhau ở đại lượng đo và ở nhóm can thiệp}. Chúng chỉ khác
nhau ở một chỗ: xử lý các biến kiểm soát thế nào.

Làm bốn cách không phải để chọn ra cái đúng nhất, mà để xem kết luận có đổi khi đổi
cách làm không.
\end{noi}

\slidemuc{15}{Cách 1: hồi quy thô}{35s}{\relax}
\begin{noi}
Không trừ gì cả, lấy thẳng chênh lệch hai nhóm. Ra \textbf{âm 0,40 điểm}.

Đây không phải để tin, mà là mốc so sánh --- để biết ba cách sau kéo con số đi được
bao nhiêu.
\end{noi}

\slidemuc{16}{Cách 2: hồi quy có hiệp biến}{40s}{\relax}
\begin{noi}
Thêm các đặc điểm nền vào cùng một phương trình, cho nó hút bớt phần khác biệt sẵn
có. Ra \textbf{âm 0,27 điểm}.

Một chi tiết đáng nói: thêm biến vào mà \textbf{sai số chuẩn lại tăng}, từ 0,60 lên
0,73. Nghe như lỗi, nhưng đúng --- vì các biến kiểm soát tương quan với biến can
thiệp, tức là đa cộng tuyến. Nhóm có kiểm lại bằng công thức và nó khớp.
\end{noi}

\slidemuc{17}{Cách 3: g-computation}{45s}{\relax}
\begin{noi}
Cách này khác hẳn về tư duy. Nhóm dựng mô hình \textbf{chỉ trên nhóm đối chứng} ---
tức là học xem ``không có chính sách thì giá đổi thế nào''.

Rồi lấy đúng mô hình đó áp lên nhóm được giảm thuế, để dự đoán: nếu nhóm này không
có chính sách thì giá lẽ ra đã đổi bao nhiêu. Lấy thực tế trừ đi dự đoán đó.

Ra \textbf{âm 0,66 điểm} --- con số lớn nhất trong bốn cách.
\end{noi}

\slidemuc{18}{Cách 4: phân tầng theo mức giá}{35s}{\relax}
\begin{noi}
Chia thành 5 tầng theo mức giá nền, chỉ so những mặt hàng giá gần nhau, rồi gộp lại
theo trọng số. Ra \textbf{âm 0,26 điểm}.
\end{noi}
\chuyen{Trước khi đọc kết quả, phải nói một chuyện.}

\slidemuc{19}{Kiểm tra thiết kế}{35s}{\relax}
\begin{noi}
Bốn cách này \textbf{không kiểm chứng lẫn nhau}, vì cả bốn cùng dựa trên một giả
định gốc là xu hướng song song. Bốn cách cùng ra kết quả giống nhau không chứng
minh giả định đó đúng.

Và nhóm phải nói thẳng: \textbf{cổng kiểm tra cân bằng đã trượt.} Đó là hạn chế
thật, lát nữa quay lại.
\end{noi}

% =====================================================================
\section*{Phần 4 --- Kết quả}

\slidemuc{20}{Bốn ước lượng}{60s}{KHÔNG CẮT}
\begin{noi}
Bốn con số: \textbf{âm 0,40 \quad âm 0,27 \quad âm 0,66 \quad âm 0,26}.

Hai điều đọc được. Thứ nhất, \textbf{cả bốn đều âm} --- tức đều cùng chiều với việc
chính sách có tác dụng phần nào.

Thứ hai, và quan trọng hơn: \textbf{cả bốn khoảng tin cậy đều phủ qua 0}, p-value
từ 0,40 tới 0,71. Nghĩa là với cỡ mẫu này, nhóm \textbf{không thể loại trừ khả năng
tác động bằng 0}.
\end{noi}
\chuyen{Từ chênh lệch này ra tỉ lệ chuyển thuế thế nào.}

\slidemuc{21}{Từ chênh lệch ra tỉ lệ}{60s}{KHÔNG CẮT --- CHỖ DỄ HIỂU SAI NHẤT}
\chidan{Nói chậm slide này. Đây là chỗ người nghe hay hiểu ngược.}
\begin{noi}
Gốc so sánh \textbf{không phải số 0}. Nếu không có chính sách, giá lẽ ra đã tăng
1,02 điểm. Thực tế nhóm được giảm thuế chỉ tăng 0,62. Vậy chính sách kéo lại được
0,40.

Còn nếu chuyển hết phần thuế vào giá thì lẽ ra phải kéo được \textbf{1,835 điểm}
--- đó là con số số học từ 1,08 chia 1,10.

Lấy 0,40 chia 1,835 ra \textbf{22\%}.

\textbf{Xin nhấn mạnh: 22\% không có nghĩa là giá giảm 22\%.} Giá vẫn tăng. 22\% là
phần trăm \emph{của mức lẽ ra kéo được}. Nói cách khác: lẽ ra giá phải thấp hơn
1,8\%, nhóm chỉ đo được 0,4\%.

Bốn cách cho khoảng \textbf{14\% đến 36\%}.
\end{noi}
\chuyen{Nhưng con số này thừa hưởng toàn bộ bất định của giả định. Nên nhóm làm
thêm một kiểm chứng không cần giả định nào.}

\slidemuc{22}{Hai câu hỏi khác nhau}{45s}{\relax}
\begin{noi}
Bảng này để tránh hiểu nhầm là nhóm đang mâu thuẫn với chính mình.

Tỉ lệ chuyển thuế là ước lượng \textbf{nhân quả} --- mạnh về ý nghĩa, nhưng phụ
thuộc giả định xu hướng song song, mà giả định đó đang lung lay.

Bám chuẩn cơ học thì \textbf{không cần giả định nào}, nhưng nó chỉ \textbf{mô tả}
--- nó không nói được nguyên nhân.

Mỗi cái mạnh đúng chỗ cái kia yếu.
\end{noi}

\slidemuc{23}{Chuẩn giá cơ học}{55s}{KHÔNG CẮT}
\begin{noi}
Cách làm rất đơn giản, không dùng nhóm đối chứng gì cả. Với mỗi mặt hàng, lấy giá
trước nhân 1,08 chia 1,10, làm tròn về bội số 1000 đồng. Đó là \textbf{giá lẽ ra
phải có}. Rồi đối chiếu với giá thật.

Trong 155 mặt hàng được giảm thuế, có \textbf{135 mặt hàng lẽ ra phải đổi giá}. Kết
quả: \textbf{đúng 1 mặt hàng} rơi vào mức đó. Một trên một trăm ba mươi lăm.

Và \textbf{110 mặt hàng giữ nguyên y nguyên giá cũ}, không xê dịch một đồng.

Nhóm có chạy giả dược trên nhóm đối chứng: 1 trên 92 --- tức là tỉ lệ bám chuẩn ở
nhóm được giảm thuế \textbf{không hề cao hơn} nhóm lẽ ra không chịu tác động nào.
\end{noi}
\chuyen{Để thấy rõ, nhóm demo trực tiếp.}

% =====================================================================
\section*{Demo trực tiếp --- 2:30}

\chidan{Mở tab \texttt{/demo}. Ba thao tác, không nói gì về phương pháp --- chỉ cho
thấy kết quả.}

\slidemuc{D1}{Món duy nhất đạt chuẩn}{50s}{\relax}
\chidan{Bấm nút gợi ý SUNLIGHT NLS Hương Lavender 997ml.}
\begin{noi}
Đây là mặt hàng duy nhất trong 135 món rơi đúng mức: 39.000 xuống 38.000.
\end{noi}

\slidemuc{D2}{Món giữ nguyên giá}{50s}{\relax}
\chidan{Bấm VASELINE Son dưỡng môi.}
\begin{noi}
Còn đây là trường hợp phổ biến hơn nhiều. Giá lẽ ra phải là 94.000. Thực tế vẫn
96.000 --- y hệt trước ngày chính sách.
\end{noi}

\slidemuc{D3}{Máy tính giá}{50s}{\relax}
\chidan{Gõ 10.000 vào ô giá.}
\begin{noi}
Và đây là một phần lý do. Với hàng giá thấp, sau khi làm tròn về bội số 1000 đồng
thì giá quay về đúng mức cũ --- mức giảm không đủ để vượt một bước làm tròn.

Nhóm có mô phỏng: nếu cửa hàng làm tròn tới 100 đồng thì cả 155 mặt hàng đều phải
đổi giá. Nên làm tròn giải thích được một phần, \textbf{nhưng không giải thích được
110 mặt hàng đứng yên}.
\end{noi}

\chidan{Nếu mạng hỏng hoặc web không lên: bỏ qua demo, nói thẳng ba con số 1 trên
135, 110 mặt hàng, và ví dụ VASELINE. Slide 23 đã có đủ.}

% =====================================================================
\section*{Phần 5 --- Hạn chế và kết luận}

\slidemuc{24}{Hạn chế}{40s}{\relax}
\begin{noi}
Ba hạn chế, nhóm nói thẳng.

Một, \textbf{cỡ mẫu nhỏ} --- 287 mặt hàng, nên khoảng tin cậy rộng và dữ liệu chưa
đủ mạnh để khẳng định có tác động.

Hai, \textbf{cổng cân bằng đã trượt}, nên giả định xu hướng song song không được dữ
liệu ủng hộ hoàn toàn.

Ba, \textbf{một cửa hàng}, không suy rộng ra cả thị trường được.

Ba hạn chế này ảnh hưởng tới \textbf{độ lớn} con số ước lượng. Chúng \textbf{không}
lật ngược được phần đếm cơ học, vì phần đó không dùng giả định nào.
\end{noi}

\slidemuc{25}{Kết luận}{35s}{\relax}
\begin{noi}
Kết luận của nhóm: \textbf{cửa hàng đã không chuyển hết phần giảm thuế vào giá bán
lẻ.}

Chỗ dựa chắc nhất không phải con số 14 đến 36\% --- con số đó còn bất định. Chỗ dựa
chắc nhất là phép đếm: \textbf{1 trên 135 mặt hàng rơi đúng mức, và 110 mặt hàng
giữ nguyên giá cũ.} Phép đếm đó không cần giả định nào.

Và dù lý do là gì đi nữa --- làm tròn, chi phí đổi nhãn, hay cửa hàng giữ lại phần
chênh --- thì \textbf{người mua vẫn đang trả giá cũ}.

Nhóm xin hết. Cảm ơn thầy và các bạn.
\end{noi}

% =====================================================================
\newpage
\section*{Hỏi đáp dự phòng}

\textbf{``Khoảng tin cậy phủ 0 mà sao vẫn kết luận được?''}
\begin{noi}
Nhóm không kết luận từ phần ước lượng nhân quả. Phần đó nhóm báo cáo là \emph{chưa
đủ bằng chứng}. Kết luận dựa trên phép đếm cơ học: 1 trên 135. Phép đếm đó không
cần giả định xu hướng song song, cũng không cần nhóm đối chứng.
\end{noi}

\textbf{``Sao không hồi quy theo thuế suất cửa hàng thực áp?''}
\begin{noi}
Vì nó là biến trung gian, nằm giữa việc đủ điều kiện giảm thuế và giá. Kiểm soát
biến trung gian sẽ chặn mất chính con đường cần đo. Nhóm ước lượng theo đủ điều
kiện pháp lý, tức là ITT.
\end{noi}

\textbf{``Bốn cách ra bốn số khác nhau, tin cái nào?''}
\begin{noi}
Không tin riêng cái nào. Khoảng 14 đến 36\% chính là câu trả lời trung thực --- nó
cho thấy kết quả nhạy với cách xử lý biến kiểm soát tới mức nào. Nếu chỉ báo một
con số thì là giấu bớt thông tin.
\end{noi}

\textbf{``Có phải chỉ vì làm tròn 1000 đồng không?''}
\begin{noi}
Một phần, với hàng giá thấp. Nhưng nhóm mô phỏng lại: ở bước 1000 đồng vẫn có 135
trên 155 mặt hàng lẽ ra phải đổi giá. Làm tròn không giải thích được 110 mặt hàng
đứng yên.
\end{noi}

\textbf{``Một cửa hàng thì nói lên được gì?''}
\begin{noi}
Không suy rộng ra thị trường được, và nhóm ghi rõ điều đó trong phần hạn chế. Nhưng
nó đủ để bác bỏ mệnh đề ``giảm thuế thì đương nhiên giá giảm'' --- chỉ cần một phản
ví dụ là đủ.
\end{noi}

\textbf{``Nhóm đối chứng là rượu bia thuốc lá, so vậy có hợp lý không?''}
\begin{noi}
Đây là điểm yếu nhóm tự nhận. Nhóm có chạy thêm ba định nghĩa nhóm đối chứng khác
nhau, kết quả vẫn cùng chiều. Nhưng cổng cân bằng vẫn trượt, nên nhóm không nhận
phần ước lượng nhân quả là bằng chứng mạnh.
\end{noi}

\textbf{``Sao thêm biến kiểm soát mà sai số lại tăng?''}
\begin{noi}
Do đa cộng tuyến --- các biến kiểm soát tương quan với biến can thiệp. Nhóm kiểm
lại bằng công thức hệ số phóng đại phương sai và con số khớp.
\end{noi}

\end{document}
